\documentclass[12pt,a4paper]{article}

% ---------- Packages ----------
\usepackage[a4paper,margin=2.5cm]{geometry}
\usepackage{setspace}
\usepackage{longtable}
\usepackage{array}
\usepackage{tabularx}
\usepackage{booktabs}
\usepackage{enumitem}
\usepackage{hyperref}
\usepackage{xcolor}
\usepackage{fancyhdr}
\usepackage{titlesec}
\usepackage{xepersian}

% ---------- Fonts ----------
\settextfont{Vazir.ttf}[
BoldFont = Vazir-Bold.ttf,
Path = fonts/
]

% ---------- Page settings ----------
\onehalfspacing
\setlength{\parindent}{0pt}
\setlength{\parskip}{6pt}
\hypersetup{
	colorlinks=true,
	linkcolor=blue,
	urlcolor=blue
}

\pagestyle{fancy}
\fancyhf{}
\fancyhead[R]{سند برنامه مدیریت پروژه نرم‌افزاری}
\fancyhead[L]{نسخه ۰.۱}
\fancyfoot[C]{\thepage}

\renewcommand{\arraystretch}{1.5}

% ---------- Document ----------
\begin{document}
	
	% ==========================================
	% Title Page
	% ==========================================
	
	\begin{titlepage}
		\centering
		
		\vspace*{2cm}
		
		{\Large سند برنامه مدیریت پروژه نرم‌افزاری}
		
		\vspace{1cm}
		
		{\huge\bfseries
			سامانه هوشمند مدیریت و خودکارسازی
			عملیات حمل‌ونقل و لجستیک
			مبتنی بر هوش مصنوعی
		}
		
		\vspace{1.5cm}
		
		{\Large $Software\ Project\ Management\ Plan\ (SPMP)$}
		
		\vspace{2cm}
		
		\begin{tabular}{rl}
			\textbf{نسخه سند:} & ۰.۱ (پیش‌نویس اولیه) \\
			\textbf{وضعیت سند:} & در حال توسعه و بازبینی \\
			\textbf{تاریخ تهیه:} & ۱۴۰۵/۰۸/۱۷ \\
			\textbf{مدل توسعه:} & افزایشی ($Incremental$) \\
		\end{tabular}
		
		\textbf{اعضای تیم:} علیرضا باجوری، علی نادرخوی ، امیرحسین پیرنژاد ، یحیی محمدزاده ، صبا محمدیان ، امین رضا امان
		\vfill
		
		
	\end{titlepage}
	
	\tableofcontents
	\newpage
	
	% ==========================================
	% Document Control
	% ==========================================

	
	\newpage
	
	% ==========================================
	% 1. Introduction
	% ==========================================
	
	\section{معرفی پروژه ($Introduction$)}
	
	\subsection{هدف سند}
	
	هدف از تهیه این سند، تعیین چارچوبی برای
	برنامه‌ریزی، سازمان‌دهی، اجرا، نظارت و کنترل
	پروژه سامانه هوشمند مدیریت و خودکارسازی
	عملیات حمل‌ونقل و لجستیک مبتنی بر هوش مصنوعی است.
	
	این سند مسئولیت اعضای تیم، فرایند توسعه،
	برنامه زمانی اولیه، منابع، ریسک‌ها و روش‌های
	کنترل کیفیت پروژه را مشخص می‌کند.
	
	همچنین این سند به‌عنوان مرجع مشترک اعضای تیم
	برای هماهنگی فعالیت‌ها و ارزیابی پیشرفت پروژه
	مورد استفاده قرار می‌گیرد.
	
	\subsection{معرفی پروژه}
	
	پروژه حاضر با هدف طراحی و توسعه یک سامانه
	نرم‌افزاری هوشمند برای مدیریت و خودکارسازی
	فرایندهای کلیدی حمل‌ونقل و لجستیک تعریف شده است.
	
	مسئله اصلی، ناکارآمدی عملیاتی، اتلاف منابع،
	هزینه‌های بالا و دشواری برنامه‌ریزی و هماهنگی
	در زنجیره حمل‌ونقل و لجستیک است.
	
	سامانه پیشنهادی با استفاده از تحلیل داده‌ها،
	الگوریتم‌های بهینه‌سازی و روش‌های هوش مصنوعی
	درصدد بهبود تخصیص منابع، برنامه‌ریزی مسیر،
	زمان‌بندی عملیات و پایش عملکرد ناوگان است.
	
	\subsection{اهداف پروژه}
	
	اهداف اصلی پروژه عبارت‌اند از:
	
	\begin{itemize}
		\item کاهش هزینه‌های عملیاتی حمل‌ونقل.
		\item بهبود تخصیص بار و وسایل نقلیه.
		\item بهینه‌سازی مسیرهای حمل‌ونقل.
		\item کاهش ترددهای بدون بار.
		\item افزایش دقت برنامه‌ریزی و زمان‌بندی.
		\item فراهم‌کردن امکان پایش و ردیابی عملیات.
		\item استفاده از تحلیل داده و هوش مصنوعی
		برای پشتیبانی از تصمیم‌گیری.
		\item ارائه گزارش‌ها و داشبوردهای مدیریتی.
	\end{itemize}
	
	\subsection{محدوده پروژه ($Project Scope$)}
	
	محدوده اولیه پروژه شامل بررسی، طراحی و توسعه
	قابلیت‌های زیر است:
	
	\begin{itemize}
		\item مدیریت کاربران و سطوح دسترسی.
		\item مدیریت ناوگان حمل‌ونقل.
		\item ثبت و مدیریت سفارش‌ها و محموله‌ها.
		\item تخصیص بار و وسیله نقلیه.
		\item مسیریابی و زمان‌بندی عملیات.
		\item پایش وضعیت عملیات و ناوگان.
		\item تحلیل داده‌ها و قابلیت‌های هوشمند.
		\item مدیریت هزینه‌ها و منابع.
		\item گزارش‌گیری و داشبورد مدیریتی.
	\end{itemize}
	
	اولویت‌بندی قابلیت‌ها و تعیین محدوده دقیق
	نسخه اولیه قابل ارائه ($MVP$) در مراحل بعدی
	و پس از تحلیل دقیق‌تر نیازمندی‌ها انجام خواهد شد.
	
	\subsection{موارد خارج از محدوده اولیه}
	
	موارد زیر در این مرحله تعهد قطعی پروژه
	محسوب نمی‌شوند و بررسی آن‌ها به تصمیم‌های
	بعدی تیم وابسته است:
	
	\begin{itemize}
		\item اتصال عملیاتی به سامانه‌های دولتی،
		گمرکی و بندری.
		\item استفاده از داده‌های زنده و تجاری
		شرکت‌های حمل‌ونقل.
		\item ارائه خدمات مالی و پرداخت واقعی.
		\item توسعه اپلیکیشن مستقل برای رانندگان.
		\item استقرار گسترده در زیرساخت‌های عملیاتی.
	\end{itemize}
	
	\subsection{ذی‌نفعان پروژه}
	
	ذی‌نفعان شناسایی‌شده عبارت‌اند از:
	
	\begin{itemize}
		\item صاحبان بار و شرکت‌های تولیدی و بازرگانی.
		\item رانندگان.
		\item شرکت‌های حمل‌ونقل.
		\item سازمان راهداری و حمل‌ونقل جاده‌ای.
		\item وزارت صنعت، معدن و تجارت.
		\item گمرک و سازمان بنادر و دریانوردی.
		\item شرکت‌های بیمه بار و ناوگان.
		\item انجمن‌های صنفی رانندگان.
	\end{itemize}
	
	\subsection{خروجی‌های پروژه ($Deliverables$)}
	
	خروجی‌های مورد انتظار پروژه شامل موارد زیر است:
	
	\begin{enumerate}
		\item مستندات تحلیل نیازمندی‌ها.
		\item مستند معماری و طراحی نرم‌افزار.
		\item طراحی پایگاه داده.
		\item کد منبع سامانه.
		\item ماژول‌های مدیریتی و عملیاتی.
		\item ماژول‌های هوشمند و بهینه‌سازی.
		\item مجموعه تست‌ها و نتایج ارزیابی.
		\item مستندات فنی و راهنمای استفاده.
		\item نسخه قابل اجرا و قابل ارائه سامانه.
	\end{enumerate}
	
	\subsection{مراجع و تعاریف}
	
	مراجع و اصطلاحات اولیه:
	
	\begin{itemize}
		\item $SPMP$: برنامه مدیریت پروژه نرم‌افزاری.
		\item $SRS$: مشخصات نیازمندی‌های نرم‌افزار.
		\item $MVP$: حداقل محصول پذیرفتنی.
		\item $AI$: هوش مصنوعی.
		\item $API$: رابط برنامه‌نویسی کاربردی.
		\item $CI/CD$: یکپارچه‌سازی و تحویل مداوم.
		\item $Git$: سیستم کنترل نسخه.
	\end{itemize}
	
	% ==========================================
	% 2. Project Organization
	% ==========================================
	
	\section{سازمان‌دهی پروژه ($Project\ Organization$)}
	
	\subsection{مدل سازمان‌دهی تیم}
	
	تقسیم دقیق مسئولیت‌ها و نقش‌های تخصصی
	در مراحل بعدی و با توجه به توانایی‌ها،
	علاقه‌مندی‌ها و حجم فعالیت‌ها تعیین خواهد شد.
	
	\subsection{نقش‌ها و مسئولیت‌ها}
	باید تکمیل شود
	
	
	\subsection{ارتباطات تیمی}
	
	ارتباطات و هماهنگی اعضا از طریق ابزارهای
	ارتباطی مورد توافق تیم انجام خواهد شد.
	
	برنامه ارتباطی اولیه:
	
	\begin{itemize}
		\item برگزاری جلسه هفتگی برای بررسی پیشرفت.
		\item ثبت تصمیم‌های مهم و تغییرات پروژه.
		\item گزارش مشکلات و موانع توسعه.
		\item مشخص‌کردن وظایف دوره بعد.
	\end{itemize}
	
	ابزارهای ارتباطی نهایی تعیین‌نشده هستند.
	
	% ==========================================
	% 3. Managerial Process
	% ==========================================
	
	\section{فرایندهای مدیریتی ($Managerial\ Process$)}
	
	\subsection{اهداف و اولویت‌های مدیریتی}
	
	اولویت‌های اصلی مدیریت پروژه عبارت‌اند از:
	
	\begin{enumerate}
		\item تکمیل نیازمندی‌ها و تعیین محدوده $MVP$.
		\item تقسیم مناسب وظایف میان اعضا.
		\item جلوگیری از گسترش کنترل‌نشده محدوده پروژه.
		\item مدیریت زمان و منابع موجود.
		\item کنترل کیفیت خروجی هر مرحله.
		\item شناسایی و کاهش ریسک‌ها.
		\item حفظ هماهنگی و مستندسازی تصمیم‌ها.
	\end{enumerate}
	
	\subsection{فرضیات پروژه}
	
	فرضیات اولیه عبارت‌اند از:
	
	\begin{itemize}
		\item اعضای تیم در طول توسعه همکاری خواهند داشت.
		\item امکان استفاده از ابزارهای توسعه نرم‌افزار
		و کنترل نسخه وجود دارد.
		\item در صورت نبود داده واقعی، امکان تولید
		یا استفاده از داده‌های آزمایشی بررسی خواهد شد.
		\item توسعه پروژه به‌صورت مرحله‌ای انجام می‌شود.
	\end{itemize}
	
	\subsection{وابستگی‌ها}
	
	پروژه به موارد زیر وابستگی دارد:
	
	\begin{itemize}
		\item دسترسی به داده‌های عملیاتی حمل‌ونقل.
		\item زیرساخت‌های مکانی و مسیریابی.
		\item زیرساخت نرم‌افزاری و سخت‌افزاری.
		\item سرویس‌ها و $API$های خارجی مورد نیاز.
		\item همکاری کاربران و دسترسی به بازخورد آن‌ها.
	\end{itemize}
	
	\subsection{محدودیت‌ها}
	
	محدودیت‌های اولیه عبارت‌اند از:
	
	\begin{itemize}
		\item محدودیت زمان توسعه.
		\item محدودیت منابع مالی و سخت‌افزاری.
		\item محدودیت دسترسی به داده‌های واقعی.
		\item پیچیدگی الگوریتم‌های هوش مصنوعی.
		\item محدودیت‌های امنیتی و قانونی.
		\item محدودیت تجربه و نیروی انسانی تیم.
	\end{itemize}
	
	\subsection{مدیریت ریسک ($Risk\ Management$)}
	
	ریسک‌های شناسایی‌شده اولیه:
	
	\begin{longtable}{|p{3.3cm}|p{4cm}|p{5.7cm}|}
		\hline
		\textbf{ریسک} & \textbf{پیامد} & \textbf{راهکار پیشنهادی} \\
		\hline
		نبود داده واقعی &
		کاهش دقت و قابلیت ارزیابی مدل‌ها &
		بررسی داده‌های عمومی و تولید داده آزمایشی \\
		\hline
		پیچیدگی الگوریتم‌ها &
		افزایش زمان توسعه &
		شروع با الگوریتم‌های ساده و توسعه تدریجی \\
		\hline
		تغییر نیازمندی‌ها &
		افزایش محدوده و تأخیر پروژه &
		ثبت تغییرات و بررسی اثر آن‌ها \\
		\hline
		عدم هماهنگی اعضا &
		تداخل وظایف و کاهش بهره‌وری &
		جلسات منظم و تعیین مسئولیت‌ها \\
		\hline
		مشکلات زیرساختی &
		اختلال در توسعه و استقرار &
		انتخاب تدریجی ابزارها و زیرساخت‌های مناسب \\
		\hline
	\end{longtable}
	
	احتمال وقوع، شدت اثر و اولویت هر ریسک
	در بازبینی‌های بعدی ارزیابی و ثبت خواهد شد.
	
	\subsection{نظارت و کنترل پروژه}
	
	نظارت بر پروژه به‌صورت هفتگی انجام خواهد شد.
	
	معیارهای اولیه ارزیابی پیشرفت عبارت‌اند از:
	
	\begin{itemize}
		\item درصد تکمیل وظایف برنامه‌ریزی‌شده.
		\item تعداد وظایف تکمیل‌شده و باقی‌مانده.
		\item میزان انطباق با برنامه زمانی.
		\item تعداد خطاهای شناسایی‌شده و رفع‌شده.
		\item وضعیت ریسک‌های مهم.
		\item میزان تکمیل مستندات.
	\end{itemize}
	
	در پایان هر دوره، وضعیت پروژه بررسی شده
	و برنامه دوره بعدی به‌روزرسانی خواهد شد.
	
	\subsection{مدیریت تغییرات}
	
	هرگونه تغییر مهم در محدوده، نیازمندی‌ها،
	معماری یا برنامه زمانی باید ثبت شود.
	
	فرایند پیشنهادی مدیریت تغییر:
	
	\begin{enumerate}
		\item ثبت درخواست تغییر.
		\item بررسی ضرورت و دلیل تغییر.
		\item ارزیابی اثر تغییر بر زمان، منابع و کیفیت.
		\item تصمیم‌گیری و توافق اعضای تیم.
		\item اعمال تغییر و به‌روزرسانی مستندات.
	\end{enumerate}
	
	% ==========================================
	% 4. Technical Process
	% ==========================================
	
	\section{فرایندهای فنی ($Technical\ Process$)}
	
	\subsection{مدل توسعه نرم‌افزار}
	
	مدل توسعه پیشنهادی پروژه،
	مدل افزایشی ($Incremental\ Model$) است.
	
	با توجه به گستردگی قابلیت‌های سامانه
	و نیاز به ارزیابی تدریجی الگوریتم‌های هوش مصنوعی،
	توسعه محصول در چند مرحله انجام خواهد شد.
	
	در هر مرحله، مجموعه‌ای از قابلیت‌ها
	طراحی، پیاده‌سازی، آزمایش و ارزیابی می‌شوند.
	
	\subsection{مراحل کلی توسعه}
	
	مراحل اولیه توسعه عبارت‌اند از:
	
	\begin{enumerate}
		\item تحلیل و تکمیل نیازمندی‌ها.
		\item طراحی معماری و پایگاه داده.
		\item توسعه قابلیت‌های پایه.
		\item توسعه قابلیت‌های مدیریتی و عملیاتی.
		\item توسعه قابلیت‌های هوشمند.
		\item یکپارچه‌سازی و تست.
		\item ارزیابی و آماده‌سازی نسخه نهایی.
	\end{enumerate}
	
	\subsection{روش‌ها و ابزارهای فنی}
	
	در ادامه تکمیل خواهد شد
	
	\subsection{ملاحظات فنی هوش مصنوعی}
	
	در ادامه تکمیل خواهد شد
	% ==========================================
	% 5. Supporting Processes
	% ==========================================
	
	\section{فرایندهای پشتیبان ($Supporting\ Processes$)}
	
	\subsection{مدیریت پیکربندی و نسخه‌ها}
	
	مدیریت نسخه‌های کد و فایل‌های پروژه
	با استفاده از سیستم کنترل نسخه انجام خواهد شد.
	
	سیاست اولیه:
	
	\begin{itemize}
		\item ثبت تغییرات در مخزن پروژه.
		\item استفاده از پیام‌های معنادار برای $Commit$ها.
		\item جلوگیری از ثبت اطلاعات محرمانه در مخزن.
		\item بازبینی تغییرات مهم پیش از ادغام.
		\item نگهداری نسخه‌های پایدار پروژه.
	\end{itemize}
	
	\subsection{تضمین کیفیت ($Quality\ Assurance$)}
	
	برای حفظ کیفیت محصول، اقدامات زیر
	به‌صورت تدریجی اجرا خواهند شد:
	
	\begin{itemize}
		\item رعایت اصول کدنویسی خوانا و قابل نگهداری.
		\item بازبینی کدهای مهم.
		\item بررسی نیازمندی‌ها پیش از پیاده‌سازی.
		\item اجرای تست برای قابلیت‌های توسعه‌یافته.
		\item ثبت و پیگیری خطاها.
	\end{itemize}
	
	\subsection{برنامه آزمون و اعتبارسنجی}
	
	آزمون‌های مورد بررسی شامل موارد زیر هستند:
	
	\begin{itemize}
		\item آزمون واحد ($Unit\ Testing$).
		\item آزمون یکپارچگی ($Integration\ Testing$).
		\item آزمون سیستم ($System\ Testing$).
		\item آزمون پذیرش ($Acceptance\ Testing$).
		\item ارزیابی عملکرد الگوریتم‌های هوشمند.
	\end{itemize}
	
	\subsection{مستندسازی}
	
	مستندات اصلی پروژه شامل موارد زیر خواهند بود:
	
	\begin{itemize}
		\item سند برنامه مدیریت پروژه ($SPMP$).
		\item سند نیازمندی‌های نرم‌افزار ($SRS$).
	\end{itemize}
	
	مستندات متناسب با پیشرفت پروژه
	به‌روزرسانی خواهند شد.
	
	% ==========================================
	% 6. Work Packages and Schedule
	% ==========================================
	
	\section{بسته‌های کاری و زمان‌بندی پروژه}
	
	\subsection{ساختار شکست کار ($WBS$)}
	
	در ادامه تکمیل خواهد شد
	
	\subsection{زمان‌بندی اولیه}
	
	زمان دقیق فعالیت‌ها در این نسخه مشخص نشده
	و پس از تخمین تلاش مورد نیاز و تعیین
	محدوده $MVP$ تکمیل خواهد شد.
	
	\subsection{منابع مورد نیاز}
	
	منابع مورد نیاز پروژه شامل موارد زیر هستند:
	
	\begin{itemize}
		\item نیروی انسانی و تخصص اعضای تیم.
		\item تجهیزات و سیستم‌های توسعه.
		\item زیرساخت ذخیره‌سازی و پردازش داده.
		\item ابزارهای توسعه و مدیریت پروژه.
		\item داده‌های مورد نیاز برای تحلیل و آزمایش.
	\end{itemize}
	
	بودجه و منابع مالی پروژه در این مرحله
	تعیین نشده‌اند.
	
	% ==========================================
	% 7. Acceptance and Completion
	% ==========================================
	
	\section{معیارهای پذیرش و تکمیل پروژه}
	
	پروژه زمانی آماده تحویل خواهد بود که
	معیارهای مورد توافق تیم و ذی‌نفعان
	برای نسخه هدف برآورده شده باشند.
	
	معیارهای اولیه عبارت‌اند از:
	
	\begin{itemize}
		\item قابلیت‌های مصوب محدوده محصول پیاده‌سازی شده باشند.
		\item عملکرد قابلیت‌های اصلی با تست بررسی شده باشد.
		\item اجزای سامانه به‌درستی یکپارچه شده باشند.
		\item خطاهای بحرانی شناسایی‌شده رفع شده باشند.
		\item مستندات فنی و راهنمای استفاده آماده باشند.
		\item نسخه قابل اجرا برای ارائه فراهم شده باشد.
	\end{itemize}
	
	معیارهای کمی و دقیق پذیرش پس از نهایی‌شدن
	نیازمندی‌ها و محدوده $MVP$ تعریف خواهند شد.
	
	% ==========================================
	% 8. Plan Evolution
	% ==========================================
	
	\section{فرایند بازبینی و تکامل سند}
	
	این سند به‌صورت مستمر و متناسب با پیشرفت
	پروژه به‌روزرسانی خواهد شد.
	
	در هر بازبینی هفتگی، موارد زیر بررسی می‌شوند:
	
	\begin{enumerate}
		\item وضعیت فعالیت‌های انجام‌شده.
		\item وظایف باقی‌مانده و موانع موجود.
		\item تغییرات احتمالی در محدوده پروژه.
		\item ریسک‌های جدید و وضعیت ریسک‌های قبلی.
		\item تصمیم‌های فنی و مدیریتی جدید.
		\item برنامه و اولویت‌های هفته آینده.
		\item بخش‌های نیازمند اصلاح در $SPMP$.
	\end{enumerate}
	
	هر نسخه جدید باید دارای شماره نسخه،
	تاریخ انتشار و شرح تغییرات باشد.
	
	% ==========================================
	% Appendix
	% ==========================================
	
	\appendix
	
	\section{موارد نیازمند تصمیم‌گیری در نسخه‌های بعدی}
	
	موارد زیر برای تکمیل برنامه مدیریت پروژه
	نیازمند بررسی و تصمیم‌گیری تیم هستند:
	
	\begin{itemize}
		\item تعیین مدیر پروژه و مسئول نگهداری اسناد.
		\item تقسیم نهایی نقش‌ها و مسئولیت‌ها.
		\item تعیین محدوده دقیق نسخه $MVP$.
		\item انتخاب زبان‌ها و فناوری‌های توسعه.
		\item انتخاب پایگاه داده.
		\item انتخاب ابزار مدیریت وظایف.
		\item تعیین تاریخ شروع و پایان پروژه.
		\item تخمین زمان و تلاش مورد نیاز هر بسته کاری.
		\item تعیین روش تأمین و تولید داده‌های آزمایشی.
		\item تعیین معیارهای کمی ارزیابی الگوریتم‌های $AI$.
		\item تعیین معیارهای دقیق پذیرش محصول.
	\end{itemize}
	
	\vspace{1cm}
	
	\begin{center}
		\textbf{پایان نسخه اولیه سند $SPMP$}
	\end{center}
	
\end{document}